\subsection{Justificación}
El crecimiento del uso de plataformas de mensajería instantánea ha transformado la forma en que las organizaciones interactúan con sus usuarios. Actualmente, canales como WhatsApp y Telegram se han convertido en medios de comunicación fundamentales para la atención al cliente, soporte técnico y gestión de consultas informativas. Sin embargo, muchas instituciones y negocios continúan gestionando estas interacciones de manera manual, lo que provoca retrasos en la respuesta, sobrecarga del personal y dificultades para mantener la consistencia en la información proporcionada.

En este contexto, el desarrollo de una plataforma basada en Software como Servicio (SaaS) que permita automatizar la atención mediante flujos de respuesta representa una solución tecnológica viable. Este tipo de sistemas permite responder de forma inmediata a preguntas frecuentes, reducir la carga operativa del personal y mejorar la disponibilidad del servicio, incluso fuera del horario laboral.

La integración de múltiples canales de mensajería dentro de una misma plataforma facilita la centralización de conversaciones, el seguimiento de consultas y la gestión eficiente de la información. En términos de adopción de mercado, WhatsApp mantiene una posición dominante a nivel global, con aproximadamente 3.000 millones de usuarios activos mensuales en octubre de 2025 \parencite{statista_messaging_2025}. Asimismo, Statista reporta 1,53 mil millones de usuarios que interactúan vía WhatsApp Business en 2023 y un ARPU aproximado de 0,25 dólares estadounidenses para ese mismo periodo \parencite{statista_whatsapp_business_2023}. Por su parte, Telegram alcanza aproximadamente 1.000 millones de usuarios activos y se posiciona como la tercera plataforma de mensajería más utilizada.

De acuerdo con LiveConnect, en su estudio sobre adopción de WhatsApp e Instagram en pymes de Latinoamérica \parencite{liveconnect_adopcion_2025}, a nivel mundial se reportan 200 millones de pymes que utilizan WhatsApp Business. Además, con base en un estudio de BCG y Meta en países en vías de desarrollo, se señala que el 75\% de las empresas utiliza WhatsApp durante todo el ciclo de vida del cliente, con liderazgo de mercados como México y Brasil. En Argentina, un estudio de 2023 encontró que el 73\% de emprendedores y pymes eran usuarios de WhatsApp, mientras que en Colombia la adopción alcanzó el 92,2\%.

\subsubsection{Justificación Económica}
La transformación digital constituye un pilar fundamental para el desarrollo económico de los países en la actualidad. En este contexto, el desarrollo de herramientas tecnológicas que faciliten la automatización de procesos de atención representa una contribución a la productividad y competitividad de la economía nacional.

La adopción de soluciones tecnológicas basadas en plataformas de mensajería automatizadas permite a las organizaciones aumentar su eficiencia operativa, lo que se traduce en una mejor utilización de los recursos disponibles a nivel nacional. Al reducir los tiempos de espera y optimizar la gestión de consultas, se favorece el incremento de la productividad en diversos sectores económicos que dependen de la atención al cliente.

Asimismo, el fortalecimiento de la infraestructura digital de atención contribuye a la inclusión tecnológica de pequeñas y medianas empresas, permitiendo que negocios de menor escala puedan acceder a herramientas de automatización que anteriormente estaban disponibles solo para grandes corporaciones. Esta democratización del acceso tecnológico favorece la equidad en el ecosistema empresarial nacional y promueve una distribución más equilibrada de los beneficios de la transformación digital.

De igual manera, la implementación de sistemas automatizados de atención mediante plataformas de mensajería contribuye a reducir la brecha digital entre los diferentes sectores económicos del país. Al facilitar herramientas accesibles y de bajo costo, se promueve la participación de más actores en la economía digital, lo que fortalece el tejido productivo nacional y contribuye al desarrollo económico inclusivo.
\subsubsection{Justificación Social }
El acceso rápido a la información y la atención oportuna se han convertido en aspectos importantes dentro de la comunicación digital actual. Muchas personas utilizan aplicaciones de mensajería instantánea como WhatsApp y Telegram para realizar consultas, solicitar servicios o comunicarse con diferentes organizaciones y negocios.

En este contexto, la implementación de herramientas tecnológicas que permitan mejorar la atención al usuario contribuye a facilitar la comunicación y el acceso a la información. El uso de sistemas automatizados de respuesta permite brindar atención inmediata a las consultas más frecuentes, mejorando la experiencia de los usuarios y reduciendo los tiempos de espera.

De esta manera, el desarrollo de la plataforma propuesta aporta una solución que favorece tanto a los usuarios como a los negocios o instituciones que ofrecen servicios, permitiendo una interacción más rápida, organizada y accesible mediante plataformas de mensajería ampliamente utilizadas.
\subsubsection{Justificación Técnica }

La implementación del sistema propuesto es técnicamente viable debido a la disponibilidad de herramientas de desarrollo, bibliotecas de programación y servicios de integración que permiten conectar aplicaciones con plataformas de mensajería como WhatsApp y Telegram. Estas tecnologías permiten diseñar sistemas capaces de procesar mensajes, generar respuestas automáticas y gestionar múltiples consultas de manera simultánea.

Asimismo, el desarrollo de este tipo de soluciones puede realizarse mediante el uso de arquitecturas de software modernas que facilitan la escalabilidad, el mantenimiento y la integración con otros sistemas digitales. Por lo tanto, el proyecto propuesto se fundamenta en el uso de tecnologías accesibles y ampliamente utilizadas en el ámbito del desarrollo de software, lo que garantiza su factibilidad técnica y su correcta implementación.
